\documentclass[11pt]{article}

\usepackage[utf8]{inputenc}
\usepackage[T1]{fontenc}
\usepackage{amsmath, amssymb, amsthm}
\usepackage{geometry}
\geometry{margin=1in}
\usepackage{hyperref}
\hypersetup{
  colorlinks=true,
  linkcolor=blue,
  citecolor=blue,
  urlcolor=blue,
}
\usepackage{mathtools}

\newtheorem{definition}{Definition}[section]
\newtheorem{theorem}[definition]{Theorem}
\newtheorem{lemma}[definition]{Lemma}
\newtheorem{proposition}[definition]{Proposition}
\newtheorem{conjecture}[definition]{Conjecture}
\newtheorem{remark}[definition]{Remark}

\newcommand{\Rnn}{\mathbb{R}_{\ge 0}}
\newcommand{\Pcal}{\mathcal{P}}
\newcommand{\Acal}{\mathcal{A}}
\newcommand{\Scal}{\mathcal{S}}
\newcommand{\Ical}{\mathcal{I}}
\newcommand{\Omk}{\Omega_k}
\newcommand{\Omsf}{\Omega_k^{\mathrm{sf}}}
\newcommand{\Omlf}{\Omega_k^{\mathrm{lf}}}
\newcommand{\Omcap}{\Omega_k^{\cap}}
\newcommand{\sigh}{\widehat{\sigma}}
\newcommand{\Tb}{T_b}
\newcommand{\Tbstar}{(\Tb)_{\!*}}

\title{The Block-Spin Step:\\
  Refined Type and Proof Obligations}
\author{Arthur Petron}
\date{\today}

\begin{document}
\maketitle

\begin{abstract}
We refine the bare \texttt{StepType} of Milestone~2.2 with the four conditions of \texttt{docs/context.md} \S3.3, abstract over the per-step budget functions $\varepsilon_{\mathrm{sf}}, \varepsilon_{\mathrm{lf}}, \Delta_{\cap}, \eta$.  The block-spin map $\Tb$ is defined as a morphism of configuration sites; the pushforward $\Tbstar$ on the bound-witness and summability-certificate presheaves is named; the dimensional-rescaling 2-cell $\Theta_k$ is defined as a separate typed object per Milestone~2.3's open-question recommendation.  The refined step type \texttt{StepTypeWithBudget} bundles the step function with the four conjuncts $S_1\!-\!S_4$, each a per-step inequality relating scale~$k$ to scale~$k+1$ via budget values at scale~$k$ only.  We state the proof obligation a term inhabiting this type must discharge and identify it as the joint deliverable of Milestones~3.1--3.4.  We do not in this document inhabit the type for any group, and we do not state any global (summability) claim about the budget.  This document corresponds to Milestone~2.3 of \texttt{work\_plan.md}.
\end{abstract}

\tableofcontents

\section{Introduction}
\label{sec:intro}

\texttt{formal/inductive-state.tex} (Milestone~2.2) defined the inductive state $\Ical_k$ and the scale-$k$-local predicate $\mathrm{Valid}_k = \mathrm{V1} \wedge \mathrm{V2} \wedge \mathrm{V3}$ (the option-B self-contained form ratified in the 2.2 close).  It also exposed a bare \texttt{StepType} = $\Ical_k \to \Ical_{k+1}$ as the clean 2.2$\to$2.3 handoff: the type was stated, the term was not.

The present document does \emph{two} things on top of that.  First, it refines the step type to carry the four conditions $S_1\!-\!S_4$ of \texttt{docs/context.md} \S3.3 as type-level constraints, abstract over the four budget functions $\varepsilon_{\mathrm{sf}}, \varepsilon_{\mathrm{lf}}, \Delta_{\cap}, \eta$.  Second, it states the categorical-level objects work\_plan.md L240--246 asks for: $\Tb$ as a site morphism, the pushforward $\Tbstar$ on $\Acal_k$ and $\Scal_k$, and the 2-cell $\Theta_k$ as a separate typed object.

\textbf{Scope.}  Everything in this document is either a type definition or a per-step inequality.  No term inhabiting the refined type is exhibited; no global / Phase-3 claim is stated.  See \S\ref{sec:scope} for the explicit boundary.

\section{The block-spin map as a site morphism}
\label{sec:site-morphism}

We work over the configuration site of \texttt{formal/sheaf-framework.tex} \S2.  At each RG scale~$k$ we have a site $\mathrm{Site}(\Omk)$ whose objects are $(U, X)$ pairs and whose covering family is the three-element small/large/buffer cover \texttt{ThreeCover}.  Let $b > 1$ be the block-scale factor and $a_{k+1} = b\, a_k$.

\begin{definition}[Block-spin configuration map]
\label{def:T-b-config}
The \emph{block-spin map at scale $k$} is a measurable function
\begin{equation*}
  \Tb \colon \Omk \;\longrightarrow\; \Omega_{k+1}
\end{equation*}
that averages link variables of the scale-$k$ lattice over blocks of side~$b$ to produce link variables on the scale-$(k+1)$ lattice.  The specific averaging procedure (\textsc{Wilson}, \textsc{Migdal--Kadanoff}, or \textsc{block-spin-with-projection}) is a program-level choice fixed in Milestone~3.x; here $\Tb$ is an abstract parameter.
\end{definition}

\begin{definition}[Block-spin polymer map]
\label{def:T-b-poly}
Each polymer $X \subseteq \Pcal_k$ has an image $\Tb(X) \subseteq \Pcal_{k+1}$ given by taking the coarse cells that intersect $X$ as the plaquette set of $\Tb(X)$.  Polymer inclusion $X \subseteq X'$ at scale $k$ induces $\Tb(X) \subseteq \Tb(X')$ at scale $k+1$.
\end{definition}

\begin{definition}[The site morphism]
\label{def:site-morphism}
The \emph{block-spin site morphism} is the functor
\begin{equation*}
  \Tb \colon \mathrm{Site}(\Omk)^{\mathrm{op}} \;\longrightarrow\; \mathrm{Site}(\Omega_{k+1})^{\mathrm{op}}
\end{equation*}
that acts on objects by $(U, X) \mapsto (\Tb(U), \Tb(X))$, on morphisms by preserving inclusion (Definition~\ref{def:T-b-poly} ensures this), and on the covering family by sending the three-element cover $\{\Omsf, \Omlf, \Omcap\}$ to the corresponding cover of $\Omega_{k+1}$ at smallness parameter $\delta_{k+1}$ (whose relation to $\delta_k$ via $b$ is fixed program-level by \texttt{docs/outline.md} \S6.4).
\end{definition}

\section{The pushforward on bound-witnesses and certificates}
\label{sec:pushforward}

\begin{definition}[Pushforward on $\Acal_k$]
\label{def:pushforward-A}
The pushforward $\Tbstar \colon \Acal_k \to \Acal_{k+1}$ assigns to each $f \in \Acal_k(U, X)$ a section $\Tbstar(f) \in \Acal_{k+1}(\Tb(U), \Tb(X))$ satisfying
\begin{equation*}
  \Tbstar(f)(\Tb(U_*)) \;\ge\; \sup\bigl\{ f(U_*') : U_*' \in \Tb^{-1}(\Tb(U_*)) \bigr\}
  \qquad \text{(as bound on the scale-$(k+1)$ activity).}
\end{equation*}
The specific form of $\Tbstar$ depends on the activity model fixed in Milestones~3.1--3.3.
\end{definition}

\begin{definition}[Pushforward on $\Scal_k$]
\label{def:pushforward-S}
For a KP certificate $(\rho, K) \in \Scal_k(U, X)$, the pushforward $\Tbstar(\rho, K) = (\rho', K') \in \Scal_{k+1}(\Tb(U), \Tb(X))$ is a new pair with $\rho'$, $K'$ related to $\rho$, $K$ by the budget conditions $S_1\!-\!S_3$ of \S\ref{sec:refined-type}.
\end{definition}

\section{The 2-cell $\Theta_k$ as a separate typed object}
\label{sec:two-cell}

Per \texttt{work\_plan.md} Milestone~2.3's open-question recommendation (``a separate typed object for clarity''), we define the dimensional-rescaling 2-cell $\Theta_k$ as its own object, not as a method on $\Ical_k$ or as an inline conjunct of the step type.

\begin{definition}[The 2-cell $\Theta_k$]
\label{def:two-cell-theta}
A \emph{$\Theta_k$-witness} is a tuple
\begin{equation*}
  \Theta_k \;=\; \bigl( \sigh_{\mathrm{in}}, \;\eta_k, \;\sigh_{\mathrm{out}}, \;\pi \bigr)
\end{equation*}
where:
\begin{itemize}
  \item $\sigh_{\mathrm{in}}, \sigh_{\mathrm{out}} \in V$ are the input and output renormalized string tensions (the $\sigma$-3 sense, \texttt{docs/context.md} \S2);
  \item $\eta_k \in V$ is the per-step degradation budget;
  \item $\pi$ is a proof of the per-step inequality
  \begin{equation}
    \label{eq:theta-decrease}
    \sigh_{\mathrm{in}} - \eta_k \;\le\; \sigh_{\mathrm{out}}.
  \end{equation}
\end{itemize}
\end{definition}

\begin{remark}[Categorical content]
$\Theta_k$ is the 2-cell of \texttt{docs/outline.md} \S10.4 expressing that ``dimensional rescaling is not a definition but a compatibility'' --- the relationship between the pushforward of the bare string tension and the next-scale renormalized one.  It carries the proof~\eqref{eq:theta-decrease} as data; its constructive role is to be inhabited by Milestone~3.4 against the concrete $\eta$ supplied there.
\end{remark}

\begin{remark}[Per-step only, no global claim]
\label{rem:theta-per-step}
Equation~\eqref{eq:theta-decrease} is a \emph{single} inequality at scale $k$.  It says nothing about $\sum_k \eta_k$, nothing about $\sigh_k$ at other scales, and nothing about a uniform lower bound $\sigh_k \ge \sigma_{\mathrm{phys}}$.  Those are Phase~3 / Milestone~3.4 claims and are explicitly out of scope here; see \S\ref{sec:scope}.
\end{remark}

\section{The refined step type}
\label{sec:refined-type}

We now refine Milestone~2.2's bare $\Ical_k \to \Ical_{k+1}$ \texttt{StepType} with the four conditions of \texttt{docs/context.md} \S3.3.  The four conditions are abstract over budget functions, which are program-level data fixed regime-by-regime in Milestones~3.1--3.4.

\begin{definition}[Per-step budget functions]
\label{def:budget-functions}
A \emph{per-step budget} at block scale $b$ is a tuple of four functions
\begin{equation*}
  (\varepsilon_{\mathrm{sf}}, \;\varepsilon_{\mathrm{lf}}, \;\Delta_{\cap}, \;\eta) \;\colon\; V \times V \;\longrightarrow\; V
\end{equation*}
each taking $(b, g_k)$ to a budget value at scale $k$.  Their concrete forms are:
\begin{itemize}
  \item $\varepsilon_{\mathrm{sf}}(b, g_k)$ --- the small-field $\rho$-degradation budget, supplied in Milestone~3.1;
  \item $\varepsilon_{\mathrm{lf}}(b, g_k)$ --- the large-field $\rho$-degradation budget, Milestone~3.2;
  \item $\Delta_{\cap}(b, g_k)$ --- the overlap $K$-growth budget, Milestone~3.3;
  \item $\eta(b, g_k)$ --- the string-tension degradation budget, Milestone~3.4.
\end{itemize}
The present document treats all four as abstract parameters.
\end{definition}

\begin{definition}[The refined step type]
\label{def:step-type-with-budget}
Fix budget functions $(\varepsilon_{\mathrm{sf}}, \varepsilon_{\mathrm{lf}}, \Delta_{\cap}, \eta)$ as in Definition~\ref{def:budget-functions}.  The \emph{block-spin step type with budget} is
\begin{equation}
  \label{eq:step-type-with-budget}
  \mathrm{StepTypeWithBudget}_b
  \;=\;
  \Bigl\{
    \mathcal{T} \colon \{\Ical_k : \mathrm{Valid}_k(\Ical_k)\} \to \{\Ical_{k+1} : \mathrm{Valid}_{k+1}(\Ical_{k+1})\}
    \;\Big|\;
    \forall \Ical_k, \;\bigl(S_1 \wedge S_2 \wedge S_3 \wedge S_4\bigr)\bigl(\Ical_k, \mathcal{T}(\Ical_k)\bigr)
  \Bigr\},
\end{equation}
where, writing $\Ical_{k+1} = \mathcal{T}(\Ical_k)$ and abbreviating $\varepsilon_*^k := \varepsilon_*(b, g_k)$, $\Delta_{\cap}^k := \Delta_{\cap}(b, g_k)$, $\eta_k := \eta(b, g_k)$:
\begin{description}
  \item[\textnormal{($S_1$) Small-field $\rho$-budget.}] $\rho_{\mathrm{sf}, k+1} \;\ge\; \rho_{\mathrm{sf}, k} - \varepsilon_{\mathrm{sf}}^k$.
  \item[\textnormal{($S_2$) Large-field $\rho$-budget.}] $\rho_{\mathrm{lf}, k+1} \;\ge\; \rho_{\mathrm{lf}, k} - \varepsilon_{\mathrm{lf}}^k$.
  \item[\textnormal{($S_3$) Overlap $K$-budget.}] $K_{\cap, k+1} \;\le\; K_{\cap, k} + \Delta_{\cap}^k$.
  \item[\textnormal{($S_4$) String-tension budget ($= \Theta_k$).}] A $\Theta_k$-witness exists with $\sigh_{\mathrm{in}} = \sigh_k$, $\eta_k = \eta(b, g_k)$, $\sigh_{\mathrm{out}} = \sigh_{k+1}$; equivalently, $\sigh_{k+1} \ge \sigh_k - \eta_k$ (Definition~\ref{def:two-cell-theta}).
\end{description}
\end{definition}

\begin{remark}[Per-step structure of each $S_i$]
\label{rem:per-step}
Each of $S_1, S_2, S_3, S_4$ is an inequality between scale-$k$ and scale-$(k+1)$ quantities, with the budget value taken at scale $k$.  No conjunct quantifies over $k$, no conjunct references $\sum_k$, no conjunct asserts a uniform-in-$k$ bound.  These are the per-step content of \texttt{docs/context.md} \S3.3; the corresponding \emph{global} claims (summability of budget; uniform lower bound on $\sigh$) are Phase~3 / Milestone~3.4, see \S\ref{sec:scope}.
\end{remark}

\begin{remark}[$\sigh_{k+1} > 0$ is not a conjunct of $S_4$]
\label{rem:no-double-positivity}
$S_4$ is the inequality $\sigh_{k+1} \ge \sigh_k - \eta_k$; it does \emph{not} additionally assert $\sigh_{k+1} > 0$.  Positivity of the output string tension is already required by the codomain of the step, which demands $\mathrm{Valid}_{k+1}(\Ical_{k+1})$; the V1 conjunct of $\mathrm{Valid}_{k+1}$ ($\sigh_{k+1} > 0$) is therefore structurally provided by the step's output type.  Asserting $\sigh_{k+1} > 0$ \emph{inside} $S_4$ would be a redundant double-statement; deriving $\sigh_{k+1} > 0$ \emph{from} $S_4$ alone would falsely conflate the per-step budget with a positivity claim that genuinely depends on the input $\sigh_k$ being large enough relative to $\eta_k$ (which is the Phase~3 budget-closure claim).
\end{remark}

\section{Proof obligations: what a term must discharge}
\label{sec:obligations}

The original P2.3 close stated the obligation atomically: a single
$\Sigma'$ bundling all four budget functions and the step.  The
P2.3-addendum decomposes this into four per-milestone
sub-obligations, with the analytic coupling exposed at the type
level.  This lets each Phase-3 milestone close its own piece
independently.

\begin{definition}[Per-milestone sub-obligations]
\label{def:sub-obligations}
The 2.3 obligation decomposes into four pieces, one per Phase-3
milestone:

\begin{description}
\item[\textnormal{\textbf{(3.1) $\mathrm{Obligation}_{\mathrm{sf}}$.}}]
Choose $\varepsilon_{\mathrm{sf}}$ and exhibit, for every valid
input $\Ical_k$, a new small-field certificate $c^{\mathrm{sf}}_{k+1}$
satisfying $S_1$.  Independent of the other three pieces.

\item[\textnormal{\textbf{(3.2) $\mathrm{Obligation}_{\mathrm{lf}}$.}}]
Choose $\varepsilon_{\mathrm{lf}}$ and exhibit $c^{\mathrm{lf}}_{k+1}$
satisfying $S_2$.  Independent of the other three pieces.

\item[\textnormal{\textbf{(3.3) $\mathrm{Obligation}_{\cap}$.}}]
Takes $\mathrm{Obligation}_{\mathrm{sf}}$ and $\mathrm{Obligation}_{\mathrm{lf}}$
as inputs.  Choose $\Delta_{\cap}$ and exhibit $c^{\cap}_{k+1}$
satisfying $S_3$ together with the V3 buffer-glue invariants
$\rho_{\cap, k+1} = \min(\rho_{\mathrm{sf}, k+1}, \rho_{\mathrm{lf}, k+1})$
and $K_{\cap, k+1} \le K_{\mathrm{sf}, k+1} + K_{\mathrm{lf}, k+1}$.
Inherits the 2.1 obligation \texttt{domGlue\_preserves\_KP}
(\texttt{lean/YMGap/Sheaf/Gluing.lean}, $\mathtt{sorryAx}$) through
V2-cap, and \texttt{sheaf-framework.tex} Conjecture~7.4 through V3-K.

\item[\textnormal{\textbf{(3.4) $\mathrm{Obligation}_{\sigma}$.}}]
Takes $\mathrm{Obligation}_{\cap}$ as input (only --- see
Remark~\ref{rem:encoded-unverified} below).  Choose $\eta$ and
exhibit $\sigh_{k+1}$ satisfying $S_4 = \Theta_k$ and the V1
positivity $\sigh_{k+1} > 0$.  Out of scope: the global closure
$\sum_k \eta(b, g_k) < \sigh_0$, which is a Phase-3-level
convergence claim about an infinite sum (see Decision~1 of the
Phase-3 kickoff for the boundary of the $\mathbb{Q}$-only
representation w.r.t.\ infinite sums).
\end{description}

The atomic obligation is recovered as the $\Sigma'$ over the four
sub-obligations:
\begin{equation*}
  \mathrm{StepObligation} \;:=\; \Sigma' \, \mathrm{sf}, \Sigma'\, \mathrm{lf}, \Sigma'\, \mathrm{cap}, \mathrm{sigma}.
\end{equation*}
A derived theorem assembling a $\mathrm{StepTypeWithBudget}_b$
inhabitant from the four pieces is left for whichever later
milestone needs it (it is not load-bearing for 3.1--3.4
individually).
\end{definition}

\begin{remark}[Encoded-but-unverified structural predictions]
\label{rem:encoded-unverified}
The dependency graph (sf $\perp$ lf; cap depends on sf, lf; sigma
depends on cap only) is a \emph{structural prediction} about
analytic derivations that do not yet exist.  It is encoded at the
Lean type level:
\begin{itemize}
\item $\mathrm{Obligation}_{\mathrm{lf}}$ takes \emph{no}
  $\mathrm{sf\_part}$ input --- the prediction is that the
  large-field budget derivation does not reference the small-field
  one.
\item $\mathrm{Obligation}_{\sigma}$ takes \emph{only}
  $\mathrm{cap\_part}$ --- the prediction is that the
  $\sigh$-transport's per-step degradation depends on the buffer
  KP constants only, not on $\varepsilon_{\mathrm{sf}}$ or
  $\varepsilon_{\mathrm{lf}}$ directly.
\end{itemize}
These predictions are \emph{self-correcting}: if 3.1/3.2 actually
couple, $\mathrm{Obligation}_{\mathrm{lf}}$ will not close at the
current signature and the type will need a coupling edge added; if
3.4 needs direct $\varepsilon$-access, $\mathrm{Obligation}_{\sigma}$'s
signature will need to be widened.  Confirmation comes when 3.2 and
3.4 close at the current signatures; refutation comes if they don't.
This remark is the audit hook for the structural prediction.
\end{remark}

\begin{remark}[A reconstructed atomic obligation]
\label{rem:atomic-obligation}
For backward reference, an inhabitant of the decomposed
$\mathrm{StepObligation}$ provides: a concrete $\Ical_{k+1}$ (built
from the four sub-obligations' outputs plus program-level data for
$g_{k+1}$, $\Pcal_{k+1}$, and $k+1$); a proof of
$\mathrm{Valid}_{k+1}(\Ical_{k+1})$ assembled from V1 (in
$\mathrm{Obligation}_{\sigma}$), V2 (each new certificate's
$\mathrm{sat}$ field), V3 (in $\mathrm{Obligation}_{\cap}$); and
proofs of $S_1, S_2, S_3, S_4$ from the corresponding
sub-obligations.  The 2.1 obligations carried through V2-cap and
V3-K are inherited as in the atomic case.
\end{remark}

\section{What this milestone does \emph{not} do}
\label{sec:scope}

To make the boundary unambiguous, we list explicitly what 2.3 \emph{does not} contain.

\begin{enumerate}
  \item \textbf{Inhabitation.}  No term of $\mathrm{StepTypeWithBudget}_b$ is exhibited, neither for the abstract value type $V$ and predicate $P$ nor for any concrete instantiation ($SU(2)$ included).  The body of the step --- constructing the next-scale certificates, proving the four budget inequalities --- is the joint content of Milestones~3.1--3.4.
  \item \textbf{Concrete budget functions.}  No specific form of $\varepsilon_{\mathrm{sf}}, \varepsilon_{\mathrm{lf}}, \Delta_{\cap}, \eta$ is supplied.  Each of these is computed regime-by-regime in Milestones~3.1, 3.2, 3.3, 3.4 respectively, from the SU(2) representation-theoretic data of Milestone~1.1 and the heat-kernel normalization of Milestone~1.2.
  \item \textbf{No global / Phase-3 budget claim.}  In particular, none of the following appears as a conjunct, assumption, or consequence anywhere in this document:
  \begin{itemize}
    \item $\sum_k \Delta_{\cap}(b, g_k) < \infty$ (the overlap-$K$-summability of \texttt{outline.md} \S10.3);
    \item $\sum_k \eta(b, g_k) < \sigh_0$ (the string-tension-budget closure of \texttt{outline.md} \S10.4);
    \item $\inf_k \sigh_k > 0$ (the uniform $\sigma$-3 survival of \texttt{outline.md} \S10.4);
    \item any statement quantified ``for all $k$'' over the RG trajectory.
  \end{itemize}
  These are Milestone~3.4's content (the inductive-budget closure).  Including any of them here would conflate per-step degradation with global summability and undo the boundary the 2.2 / 2.3 / 3.4 staging exists to preserve.
\end{enumerate}

\section{Lean mirror}
\label{sec:lean}

The Lean formalization is in \texttt{lean/YMGap/RG/BlockSpin.lean}:
\begin{itemize}
  \item \texttt{TwoCellTheta} as a \texttt{structure} carrying $(\sigh_{\mathrm{in}}, \eta, \sigh_{\mathrm{out}}, \mathrm{decrease})$ per Definition~\ref{def:two-cell-theta}.  Separate typed object per work\_plan's recommendation.
  \item \texttt{StepTypeWithBudget} as a \texttt{def} returning a \texttt{Type} (the $\Sigma'$-typed family of step functions with their S-conjuncts), abstract over the budget functions $\varepsilon_*, \Delta_{\cap}, \eta$ per Definition~\ref{def:budget-functions}.
  \item \texttt{StepObligation} as the $\Sigma'$-typed obligation: a budget and a $\mathrm{StepTypeWithBudget}$ inhabitant.  This is the explicit type that Milestone~3.x must produce.
  \item \texttt{step\_obligation\_witness} as a \texttt{sorry}'d \texttt{def} of type \texttt{StepObligation}.  The \texttt{sorry} is the unwritten 3.x term, logged in \texttt{lean/SORRY\_LEDGER.md} as \texttt{[P2.3]}, owed by Milestones~3.1--3.4.
\end{itemize}
The value type $V$ requires \texttt{[Sub V]} in addition to the 2.2 typeclasses (\texttt{[LT V] [Zero V] [LE V] [Min V] [Add V]}); subtraction is needed to express $\rho_k - \varepsilon^k_*$ and $\sigh_k - \eta_k$.  For Nat-valued instantiations \texttt{[Sub V]} is truncating; for $\mathbb{Z}$, $\mathbb{Q}$, or constructive $\mathbb{R}$ it is exact.  The refined type does not care which: each $S_i$ is a single $\le$-inequality, and the step's codomain demands $\mathrm{Valid}_{k+1}$ which ensures positivity of $\sigh_{k+1}$ independently of the sign of $\sigh_k - \eta_k$.

\section{Cross-references}
\label{sec:xref}

\begin{itemize}
  \item Prose companion: \texttt{docs/block-spin-morphism.md}.
  \item Lean formalization: \texttt{lean/YMGap/RG/BlockSpin.lean}.
  \item Python representation: \texttt{computations/block\_spin.py} with tests in \texttt{tests/test\_block\_spin.py}.
  \item Diagram: \texttt{visualizations/structure/block-spin-2cell.tikz}.
  \item Foundations: \texttt{formal/inductive-state.tex} (Definition 2.1 the state, 3.1 \texttt{Valid}, 4.1 the bare \texttt{StepType} of 2.2).
  \item Framework: \texttt{formal/sheaf-framework.tex} (Definition 5.1 compatibility, Proposition 7.1 buffer summability, Conjecture 7.4 tightness).
  \item Specification: \texttt{docs/context.md} \S3.3 (the four conditions), \S2 ($\sigma$-3 sense).  \texttt{docs/outline.md} \S10.3--\S10.4 (block-spin obligations and string-tension transport), \S10.5 (Phase-3 main theorem target).
  \item Open obligations: \texttt{lean/SORRY\_LEDGER.md} entry \texttt{[P2.3]} (this milestone, owed by 3.1--3.4), and the inherited 2.1 \texttt{[P2.1]} entry (parts (a) and (b)) which surfaces through V2-cap and V3-K of $\mathrm{Valid}_{k+1}$.
\end{itemize}

\end{document}
